% ============================================================
% 60_boxes.tex — tcolorbox-Stile, Balken, Boxen, Formelnoten, Listen
% ============================================================

\newcommand{\BreakIfNotEnoughSpace}{%
  \par\penalty0\vspace{0pt plus \ZSFbreakThreshold}\penalty-100\vspace{0pt plus -\ZSFbreakThreshold}\relax
}

% Strenger Break-Trigger für ChapterBar: grössere Mindesthöhe verhindert,
% dass ein neues Kapitel direkt am Spalten-/Seitenende beginnt.
\newcommand{\BreakIfNotEnoughSpaceChapter}{%
  \par\penalty0\vspace{0pt plus \ZSFbreakThresholdChapter}\penalty-100\vspace{0pt plus -\ZSFbreakThresholdChapter}\relax
}

% Einheitlicher Top-Spacing-Kontrakt für Titelbalken.
% Reihenfolge ist kritisch: \addvspace MUSS vor \BreakIfNotEnoughSpace stehen,
% damit es per Max-Semantik mit dem after-skip der Vorbox kollabiert.
% tcolorbox-Definitionen, die diesen Helper rufen, MUESSEN before skip=0pt
% setzen, damit das Spacing nicht doppelt eingerechnet wird.
\newcommand{\ZSFTitleBarTopSpacing}[1]{%
  \par
  \addvspace{#1}%
  \BreakIfNotEnoughSpace
}

% ------------------------------------------------------------
% Chapter-Bar
% ------------------------------------------------------------
\newtcolorbox{chapterbar}[1][]{
  enhanced jigsaw,
  breakable=false,
  colback=\ChapterBarColor,
  colframe=\ChapterBarColor,
  boxrule=0pt,
  boxsep=0pt,
  arc=\ZSFcornerBox, outer arc=\ZSFcornerBox,
  left=\ZSFspaceS,
  right=\ZSFspaceS,
  top=\ZSFspaceS,    % MaschKo-Verdichtung: S (erhöht um 1pt)
  bottom=\ZSFspaceS, % MaschKo-Verdichtung: S
  before skip=\ZSFspaceM, % MaschKo-Verdichtung: L -> M
  after skip=\ZSFchapterBarAfterSkip,
  fontupper=\ZSFfontChapter,
  #1
}
\newcommand{\ChapterBar}[1]{
  \BreakIfNotEnoughSpaceChapter
  \begin{chapterbar}
    \textcolor{\ChapterBarTextColor}{{\ZSFVariableColorsOff #1}}
  \end{chapterbar}
}

% ------------------------------------------------------------
% Dokument-Titel-Masthead — einmalig vor dem Front-Chapter
% ------------------------------------------------------------
\newtcolorbox{zsftitleheader}[1][]{
  enhanced jigsaw,
  breakable=false,
  colback=ZSFDocumentTitleBarColor,
  colframe=ZSFDocumentTitleBarColor,
  boxrule=0pt,
  boxsep=0pt,
  arc=\ZSFcornerBox, outer arc=\ZSFcornerBox,
  left=\ZSFspaceM,
  right=\ZSFspaceM,
  top=\ZSFspaceS,
  bottom=\ZSFspaceS,
  before skip=0pt,
  after skip=\ZSFspaceXS,
  halign=center,
  #1
}
\newcommand{\ZSFTitleHeader}[2]{%
  \penalty-500
  \begin{zsftitleheader}
    \sffamily
    {\color{ZSFDocumentTitleTextColor}\fontsize{11bp}{12.5bp}\bfseries\selectfont\ZSFVariableColorsOff #1\par}%
    \vspace{1pt}%
    {\color{ZSFDocumentTitleBylineColor}\fontsize{6.7bp}{7.5bp}\itshape\selectfont #2\par}%
  \end{zsftitleheader}
  \penalty10000
}

% ------------------------------------------------------------
% Subsection-Bar (+ unnummerierte Star-Variante)
% ------------------------------------------------------------
\newtcolorbox{subsectionbar}[1][]{
  enhanced jigsaw,
  breakable=false,
  colback=\SubsectionBarColor,
  colframe=\SubsectionBarColor,
  boxrule=0pt,
  boxsep=0pt,
  arc=\ZSFcornerBox, outer arc=\ZSFcornerBox,
  left=\ZSFspaceS,
  right=\ZSFspaceS,
  top=\ZSFspaceS,    % MaschKo-Verdichtung: S (erhöht um 1pt)
  bottom=\ZSFspaceS, % MaschKo-Verdichtung: S
  before skip=0pt,
  after skip=\ZSFsubsectionAfterGap,
  fontupper=\ZSFfontSection,
  #1
}
\newtcolorbox{subsubsectionbar}[1][]{
  enhanced jigsaw,
  breakable=false,
  colback=\SubsubsectionBarColor,
  colframe=\SubsubsectionBarColor,
  boxrule=0pt,
  boxsep=0pt,
  arc=\ZSFcornerBox, outer arc=\ZSFcornerBox,
  left=\ZSFspaceS,
  right=\ZSFspaceS,
  top=\ZSFspaceXS,
  bottom=\ZSFspaceXS,
  before skip=0pt,
  after skip=\ZSFsubsectionAfterGap,
  fontupper=\ZSFfontSubsubsection,
  #1
}
\makeatletter
\newcommand{\SubsectionBar}{\@ifstar{\SubsectionBarStar}{\SubsectionBarNumbered}}
\newcommand{\SubsectionBarNumbered}[2][]{%
  \ZSFTitleBarTopSpacing{\ZSFsubsectionBarBeforeSkip}
  \penalty-500
  \ZSFCollapseAfterChapterBar
  \refstepcounter{subsection}
  \setcounter{subsubsection}{0}
  \if\relax\detokenize{#1}\relax\else\label{#1}\fi
  \begin{subsectionbar}
    \textcolor{\SubsectionBarTextColor}{{\ZSFVariableColorsOff\thesubsection. #2}}
  \end{subsectionbar}
  \penalty10000
  \ZSFMarkAfterSubsectionBar
}
\newcommand{\SubsectionBarStar}[2][]{%
  \par\addvspace{\ZSFsubsectionBarBeforeSkip}%
  \penalty-200
  \ZSFCollapseAfterChapterBar
  \if\relax\detokenize{#1}\relax\else\label{#1}\fi
  {\ZSFfontSection\color{\SubsectionBarColor}\ZSFVariableColorsOff #2\par}%
  \penalty10000
  \ZSFMarkAfterSubsectionBar
}
% Variante: Abschnitt auf neuer Spalte beginnen (ersetzt rohes \columnbreak)
\newcommand{\SubsectionBarOnNewColumn}{\@ifstar{\SubsectionBarOnNewColumnStar}{\SubsectionBarOnNewColumnNumbered}}
\newcommand{\SubsectionBarOnNewColumnNumbered}[2][]{%
  \columnbreak
  \SubsectionBarNumbered[#1]{#2}%
}
\newcommand{\SubsectionBarOnNewColumnStar}[2][]{%
  \columnbreak
  \SubsectionBarStar[#1]{#2}%
}

% Kompakte dritte Ebene: nummeriert als x.x.x mit flachem Farbbalken.
\newcommand{\SubsubsectionBar}{\@ifstar{\SubsubsectionBarStar}{\SubsubsectionBarNumbered}}
\newcommand{\SubsubsectionBarNumbered}[2][]{%
  \par\addvspace{\ZSFsubsubsectionBarBeforeSkip}%
  \penalty-200
  \refstepcounter{subsubsection}%
  \if\relax\detokenize{#1}\relax\else\label{#1}\fi
  \begin{subsubsectionbar}
    \textcolor{\SubsubsectionTitleColor}{{\ZSFVariableColorsOff\thesubsubsection. #2}}
  \end{subsubsectionbar}
  \penalty10000
  \ZSFMarkAfterSubsectionBar
}
\newcommand{\SubsubsectionBarStar}[2][]{%
  \par\addvspace{\ZSFsubsubsectionBarBeforeSkip}%
  \penalty-200
  \if\relax\detokenize{#1}\relax\else\label{#1}\fi
  \begin{subsubsectionbar}
    \textcolor{\SubsubsectionTitleColor}{{\ZSFVariableColorsOff #2}}
  \end{subsubsectionbar}
  \penalty10000
  \ZSFMarkAfterSubsectionBar
}
\newcommand{\SubsubsectionBarOnNewColumn}{\@ifstar{\SubsubsectionBarOnNewColumnStar}{\SubsubsectionBarOnNewColumnNumbered}}
\newcommand{\SubsubsectionBarOnNewColumnNumbered}[2][]{%
  \columnbreak
  \SubsubsectionBarNumbered[#1]{#2}%
}
\newcommand{\SubsubsectionBarOnNewColumnStar}[2][]{%
  \columnbreak
  \SubsubsectionBarStar[#1]{#2}%
}
\makeatother

% ------------------------------------------------------------
% Gemeinsame tcb-Stile
% ------------------------------------------------------------
\tcbset{zsftitlebar/.style={
  fonttitle=\ZSFfontBoxTitle\strut\ZSFVariableColorsOff,
  halign title=center,
  title style={minimum height=\ZSFtitleBarHeight},
  lefttitle=\ZSFspaceXS,
  righttitle=\ZSFspaceXS,
  % MaschKo-Verdichtung: kein extra Titel-Padding — der \strut in fonttitle
  % trägt die Balkenhöhe bereits (spart ~2pt pro Titelbox)
  toptitle=0pt,
  bottomtitle=0pt,
  before title={},
  after title={},
}}

% Rechtsbuendiges Meta-/Kategorie-Tag im Titelbalken einer Titelbox.
% Verwendung im Titel-Argument der Box, z.B.:
%   \begin{defbox}[Quicksort\ZSFTitleTag{O(n log n)}] ... \end{defbox}
%   \begin{codebox}[Algorithmus\ZSFTitleTag{Python}]  ... \end{codebox}
% Schiebt das Tag per \hfill an den rechten Titelrand. Schrift: \ZSFfontTitleTag.
\newcommand{\ZSFTitleTag}[1]{\hfill{\ZSFfontTitleTag #1}}

\tcbset{zsfpanelbox/.style={
  enhanced jigsaw,
  colback=white,
  colbacktitle=\BoxTitleColor,
  colframe=\BoxTitleColor,
  coltitle=\BoxTitleTextColor,
  boxrule=0.5pt,
  titlerule=0pt,
  boxsep=0pt,
  arc=\ZSFcornerBox, outer arc=\ZSFcornerBox,
  before skip=\ZSFBoxBeforeSkip,
  after skip=\ZSFspaceXS, % MaschKo-Verdichtung: S -> XS
  before={\ZSFConsumeAfterSubsectionBar},
  zsftitlebar,
}}

% Single-Point-of-Truth für Titelboxen (defbox/tablebox/figbox/warnbox).
% left/right/top/bottom HIER setzen, damit \linewidth in jeder Titelbox
% identisch ist (sonst variierende Tabellen-Breiten je Box-Typ).
\tcbset{zsftitlebox/.style={
  enhanced jigsaw,
  breakable=false,
  colback=\BoxBackColor,
  colbacktitle=\BoxTitleColor,
  colframe=\BoxFrameColor,
  coltitle=\BoxTitleTextColor,
  boxrule=0.5pt,
  titlerule=0pt,
  boxsep=0pt,
  arc=\ZSFcornerBox, outer arc=\ZSFcornerBox,
  left=\ZSFspaceS,
  right=\ZSFspaceS,
  top=\ZSFspaceXS,    % MaschKo-Verdichtung: S -> XS (vertikales Chrome)
  bottom=\ZSFspaceXS, % MaschKo-Verdichtung: S -> XS
  before skip=\ZSFBoxBeforeSkip,
  after skip=\ZSFspaceXS, % MaschKo-Verdichtung: S -> XS
  before={\ZSFConsumeAfterSubsectionBar},
  zsftitlebar,
  before upper={\parskip=\ZSFspaceXS\relax\ZSFReadableAuto}, % MaschKo-Verdichtung: S -> XS (6-Seiten-Ziel)
}}

% tablebox: zsftitlebox-Rahmen/Ecken, aber ohne Innen-Padding um die Tabelle.
\tcbset{zsfTableboxShell/.style={
  zsftitlebox,
  left=0pt,
  right=0pt,
  top=0pt,
  bottom=0pt,
  colback=white,
  before upper={},
}}

% ------------------------------------------------------------
% Titelboxen
% ------------------------------------------------------------
\newtcolorbox{defbox}[1][]{zsftitlebox, title={#1}}
\newtcolorbox{tablebox}[1][]{zsfTableboxShell, title={#1}}
\newtcolorbox{figbox}[1][]{
  enhanced jigsaw,
  breakable=false,
  colback=\FigboxBackColor,
  frame hidden,
  boxrule=0pt,
  boxsep=0pt,
  arc=\ZSFcornerBox, outer arc=\ZSFcornerBox,
  left=0pt,
  right=0pt,
  top=\ZSFtextMinPad,
  bottom=\ZSFtextMinPad,
  before skip=\ZSFspaceXS,
  after skip=\ZSFspaceXS,
  before={\ZSFConsumeAfterSubsectionBar},
  before upper={%
    \centering
    \if\relax\detokenize{#1}\relax\else
      {\setlength{\fboxsep}{0pt}%
       \colorbox{\FigboxTitleBackColor}{\makebox[\linewidth][c]{\ZSFfontBoxTitle\color{\FigboxTitleTextColor}\ZSFVariableColorsOff #1}}\par}%
      \vspace{-\parskip}\vspace{\ZSFspaceXXS}%
    \fi
  },
}

% ------------------------------------------------------------
% Verfahrens-Baum (Box-Matrix): Verfahren (solide Box) + Beispiele
% (helle solide Box, linksbuendig) nebeneinander, beliebig viele Zeilen.
% Zeilenhöhen passen sich durch die TikZ-matrix automatisch an den
% Inhalt an; Spaltenbreiten sind relativ zu \linewidth (= Spaltenbreite
% im multicols-Layout) gesetzt -> kein \scalebox-Hack, volle Lesbarkeit.
% Als benannter tikzset-Stil (nicht Environment/Macro-Argument) definiert:
% \matrix braucht seine &/\\-Tokens im nativen Catcode-Kontext des
% Dokuments, das geht nur, wenn die matrix direkt im Kapitel steht.
% Nach der matrix (Name "mat") verbindet \ZSFverfahrenconnectors{<Zeilenzahl>}
% jede Verfahren-Box mit ihrer Beispiele-Box sowie alle Verfahren-Boxen
% untereinander über einen kurzen Stamm links — Anlehnung an den
% urspruenglichen Baum, aber ohne separate Gruppen-Spalte (die Kategorie
% steht bereits im \SubsubsectionBar* darüber).
% Nutzung (siehe z.B. chapters/ch10_fertigung.tex):
%   \SubsubsectionBar*{Titel}
%   \ZSFverfahrenBefore
%   {\centering
%   \begin{tikzpicture}[font=\ZSFfontTableDense]
%   \matrix[ZSFverfahrenmatrix, name=mat] {
%     Verfahren A & Beispiel A \\
%     Verfahren B & Beispiel B \\
%   };
%   \ZSFverfahrenconnectors{2}
%   \end{tikzpicture}\par}
% Mehrzeiliger Zellinhalt: \newline statt \\ verwenden (sonst Konflikt
% mit dem Zeilentrenner der matrix).
\tikzset{ZSFverfahrenmatrix/.style={
  matrix of nodes, nodes={align=center, inner sep=\ZSFspaceXS, very thin},
  row sep=\ZSFspaceS, column sep=\ZSFspaceM,
  column 1/.style={nodes={draw, rounded corners=\ZSFcornerBox, text width=0.33\linewidth}},
  column 2/.style={nodes={draw, rounded corners=\ZSFcornerBox, fill=black!4, draw=black!25, align=left, text width=0.52\linewidth}},
}}
\newcommand{\ZSFverfahrenconnectors}[1]{%
  \draw[thin,gray!60] ($(mat-1-1.west)+(-\ZSFspaceM,0)$) -- ($(mat-#1-1.west)+(-\ZSFspaceM,0)$);
  \foreach \ZSFvfrow in {1,...,#1} {
    \draw[thin,gray!60] ($(mat-\ZSFvfrow-1.west)+(-\ZSFspaceM,0)$) -- (mat-\ZSFvfrow-1.west);
    \draw[thin,gray!60] (mat-\ZSFvfrow-1.east) -- (mat-\ZSFvfrow-2.west);
  }
}
% Kollabiert den Abstand zwischen \SubsubsectionBar* und der Matrix
% (parskip-Trim, analog zu anderen Box-Titeln in diesem Modul).
\newcommand{\ZSFverfahrenBefore}{\par\vspace{-\parskip}\vspace{\ZSFspaceXXS}}

% Warn-Box (rot) — für Hinweise und Stolperfallen
\newtcolorbox{warnbox}[1][Achtung]{
  zsftitlebox,
  colback=WarnboxBack,
  colbacktitle=ETHRot,
  coltitle=white,
  title=#1,
}

% Randlose Tabelle mit Titelleiste (volle Box-Breite, kein Innenrand)
\newtcolorbox{fulltablebox}[2][]{
  zsfpanelbox,
  colframe=black!20,
  breakable=false,
  title={#1},
  boxsep=0pt, left=0pt, right=0pt, top=\ZSFtextMinPad, bottom=\ZSFtextMinPad,
  tabularx*={\arrayrulecolor{\ZSFtableRuleColor}\setlength{\arrayrulewidth}{\ZSFtableRuleWidth}}{#2},
  zsftitlebar,
  boxrule=0.5pt
}

% Randlose Box ohne Titel (volle Breite, zentriert)
\newtcolorbox{fulllessbox}[1]{
  enhanced jigsaw,
  breakable=false,
  colback=white,
  colframe=white,
  boxrule=0pt,
  arc=\ZSFcornerBox, outer arc=\ZSFcornerBox,
  before skip=\ZSFspaceS,
  after skip=\ZSFspaceXS, % MaschKo-Verdichtung: S -> XS
  boxsep=0pt, left=0pt, right=0pt, top=\ZSFtextMinPad, bottom=\ZSFtextMinPad,
  tabularx*={\arrayrulecolor{\ZSFtableRuleColor}\setlength{\arrayrulewidth}{\ZSFtableRuleWidth}}{#1},
  boxrule=0pt
}

% ------------------------------------------------------------
% Formula-Box + Longformula
% ------------------------------------------------------------
\newenvironment{longformula}{\[\begin{aligned}}{\end{aligned}\]}
% formulabox gehört zur Titelbox-Familie: optionales Argument ist ein Titel
% im zsftitlebar-Stil (wie defbox/tablebox). Der getoente Hintergrund
% (\FormulaboxBackColor) bleibt das Erkennungsmerkmal für Formelinhalt.
\tcbset{zsfformulastyle/.style={
  enhanced jigsaw,
  colback=\FormulaboxBackColor,
  frame hidden,
  boxrule=0pt,
  boxsep=0pt,
  arc=\ZSFcornerBox, outer arc=\ZSFcornerBox,
  left=\ZSFspaceXS,
  right=\ZSFspaceXS,
  top=\ZSFtextMinPad,
  bottom=\ZSFtextMinPad,
  before skip=\ZSFspaceXS,
  after skip=\ZSFspaceXS,
  before={\ZSFConsumeAfterSubsectionBar},
  before upper={\parskip=\ZSFspaceXXS\centering}, % kompakter Formelzeilen-Abstand
  after upper={\par},
  middle=\ZSFspaceSep,
}}
\newcommand{\ZSFCompactDisplayMathOn}{%
  \def\[{\par\noindent\hbox to \linewidth\bgroup\hfil$\displaystyle}%
  \def\]{$\hfil\egroup\par}% chktex 49
}
\NewDocumentEnvironment{formulabox}{O{}}
  {\begingroup
   \ZSFCompactDisplayMathOn
   \begin{tcolorbox}[zsfformulastyle]%
   \if\relax\detokenize{#1}\relax\else
     {\ZSFfontBoxTitle\color{TextVorBoxColor}\ZSFVariableColorsOff #1\par}%
     \vspace{-\parskip}\vspace{\ZSFspaceXXS}%
   \fi}
  {\end{tcolorbox}\endgroup}

% Formel-Note (klein, kursiv, dezent) — unter oder über Formeln
\newcommand{\formulanote}[1]{%
  \par\vspace{-\parskip}\vspace{\ZSFspaceXXS}% MaschKo-Verdichtung: S -> XS -> XXS
  \noindent{{\ZSFfontNote #1}}%
  \par\vspace{-\parskip}\vspace{\ZSFspaceXXS}% MaschKo-Verdichtung: S -> XS -> XXS
}
\let\formulanoteabove\formulanote
% Quellen der verwendeten Eingangsgrössen. Eine gemeinsame horizontale Zeile
% spart Platz und macht die Abhängigkeitsart via \ZSFsource explizit.
\newcommand{\formulasources}[1]{%
  \formulanoteabove{{\upshape\bfseries Eingaben:} #1}%
}

% Duenne, graue Trennlinie zwischen Formeln in einer formulabox.
% Als \hrule im vertikalen Modus gesetzt: eine Rule als Textabsatz wuerde
% eine volle Grundlinien-Höhe (~6pt) verbrauchen — so kostet sie nur
% 2x \ZSFspaceSep + 0.5pt.
\newcommand{\formulasep}{%
  \par\vspace{-\parskip}\vspace{\ZSFspaceSep}%
  {\color{black!30}\hrule height 0.5pt}%
  \vspace{\ZSFspaceSep}%
  \penalty-50%
}

% Zentriertes Item-Heading mit Trennstrich — für gleichwertige Fälle in Boxen
\newcommand{\ZSFItemHeading}[1]{%
  \par\vspace{-\parskip}\vspace{\ZSFspaceS}% MaschKo-Verdichtung: M -> S
  {\centering\textbf{#1}\par}\vspace{-\parskip}\vspace{\ZSFspaceXS}%
  \noindent\rule{\linewidth}{0.5pt}\par\vspace{-\parskip}\vspace{\ZSFspaceXS}%
}

% ------------------------------------------------------------
% Runintext + Grafik-Helpers
% ------------------------------------------------------------
\newenvironment{runintext}{%
  \ZSFConsumeAfterChapterBar
  \par\addvspace{\ZSFspaceXS}\noindent\ZSFReadableAuto\ignorespaces % MaschKo-Verdichtung: S -> XS
}{\par}

\newcommand{\PlaceholderGraphic}[1]{
  \fbox{\parbox{0.9\linewidth}{\centering\small [Grafik: #1]}}
}

% Dezente Bildunterschrift (leerer Text -> nichts)
\newcommand{\ZSFfigcaption}[1]{%
  \if\relax\detokenize{#1}\relax\else
    \par\vspace{\ZSFspaceXS}%
    {\ZSFfontNote\color{black!60}#1\par}%
  \fi
}

% \ZSFfig[breite]{pfad}{Titel}{Caption} — figbox mit Bild + optionaler Caption.
% Ohne optionales Argument: Bild auf min(Spaltenbreite, ZSFfigMaxHeight) skaliert.
% Mit [breite] (Bruchteil von \linewidth, z.B. [0.7]): exakte Breite ohne
% Höhen-Cap — erhaelt das gewünschte Grössenverhältnis der Vorlage.
\newcommand{\ZSFfig}[4][]{%
  \begin{figbox}[{#3}]%
    \centering
    \if\relax\detokenize{#1}\relax
      \includegraphics[width=\ZSFfigScale\linewidth,height=\ZSFfigScale\ZSFfigMaxHeight,keepaspectratio]{#2}%
    \else
      \includegraphics[width=\ZSFfigScale\dimexpr#1\linewidth\relax,keepaspectratio]{#2}%
    \fi
    \ZSFfigcaption{#4}%
  \end{figbox}%
}

% \ZSFfigside[Bildbreite]{pfad}{Titel}{Inhalt rechts} — Bild links, Inhalt rechts.
% Optionales [Bildbreite]-Argument (Bruchteil von \linewidth) überschreibt die
% Standardbreite von 30% -- Textspalte fuellt den Rest via \hfill.
\newlength{\ZSFfigsideImgWidth}\setlength{\ZSFfigsideImgWidth}{0.30\linewidth}
\newlength{\ZSFfigsideTextWidth}\setlength{\ZSFfigsideTextWidth}{0.66\linewidth}
\newcommand{\ZSFfigside}[4][0.30]{%
  \begin{figbox}[{#3}]%
    \noindent
    \begin{minipage}[c]{#1\linewidth}%
      \raggedright
      \includegraphics[width=\linewidth,height=\ZSFfigMaxHeight,keepaspectratio]{#2}%
    \end{minipage}%
    \hfill
    \begin{minipage}[c]{\dimexpr0.96\linewidth-#1\linewidth\relax}%
      \raggedright
      #4%
    \end{minipage}%
    \par
  \end{figbox}%
}

% Side-by-side-Layout (Bild + Text in gleicher Linie) — sehr platzsparend
\NewDocumentEnvironment{splitbox}{O{0.3} O{}}{%
  \begin{tcolorbox}[
    blankest,
    sidebyside,
    sidebyside align=top,
    sidebyside gap=\ZSFspaceM,
    lefthand width=#1\linewidth,
    lower separated=false,
    before={\ZSFConsumeAfterSubsectionBar},
    before skip=\ZSFspaceXS,
    after skip=\ZSFspaceXS,
    boxsep=0pt, left=0pt, right=0pt, top=\ZSFtextMinPad, bottom=\ZSFtextMinPad,
    #2
  ]%
}{%
  \end{tcolorbox}%
}

% ------------------------------------------------------------
% Aussagen / Verfahren / Listen
% ------------------------------------------------------------
% statementbox: dezent, linker Farbbalken in Kapitelfarbe.
\tcbset{zsfstatementstyle/.style={
  enhanced jigsaw,
  breakable=false,
  colback=white,
  colframe=white,
  frame hidden,
  borderline west={\ZSFstatementBarWidth}{0pt}{\StatementBarColor},
  boxrule=0pt,
  arc=\ZSFcornerBox, outer arc=\ZSFcornerBox,
  left=\ZSFstatementLeftPad,
  right=\ZSFspaceS,
  top=\ZSFspaceXS,
  bottom=\ZSFspaceXS,
  before skip=\ZSFBoxBeforeSkip,
  after skip=\ZSFspaceXS, % MaschKo-Verdichtung: S -> XS
  before={\ZSFConsumeAfterSubsectionBar},
}}

\NewDocumentEnvironment{statementbox}{O{}}
  {\begin{tcolorbox}[zsfstatementstyle]%
   \if\relax\detokenize{#1}\relax\else
     {\bfseries\color{\StatementTitleColor}\ZSFVariableColorsOff #1}\par\vspace{\ZSFspaceXS}%
   \fi\ignorespaces}
  {\end{tcolorbox}}

% procedure: nummerierte Schritt-Liste in einer statementbox.
\NewDocumentEnvironment{procedure}{O{}}
  {\begin{statementbox}[#1]%
   \begin{enumerate}[leftmargin=1.6em,labelsep=0.4em,itemsep=\ZSFprocItemSep,topsep=\ZSFspaceXS,parsep=0pt,partopsep=0pt]}
  {\end{enumerate}\end{statementbox}}
\newcommand{\ProcStep}[2]{\item {\bfseries #1}\\[\ZSFprocStepGap]#2}

% factlist: rahmenlose Liste von Definitions-/Eigenschaftsfakten.
\NewDocumentEnvironment{factlist}{}
  {\begin{itemize}[
     leftmargin=1.2em,
     labelsep=0.4em,
     label={\textcolor{\StatementBarColor}{\textbullet}},
     itemsep=\ZSFprocItemSep,
     topsep=\ZSFspaceXS,
     parsep=0pt,
     partopsep=0pt
   ]}
  {\end{itemize}}
\newcommand{\ZSFFact}[2]{\item \textbf{#1}\,---\,#2}

% propertylist: factlist innerhalb einer statementbox (mit optionalem Titel)
\NewDocumentEnvironment{propertylist}{O{}}
  {\par\addvspace{\ZSFspaceXS}%
   \if\relax\detokenize{#1}\relax\else
     {\bfseries\color{\StatementTitleColor}\ZSFVariableColorsOff #1}\par\vspace{-\parskip}\vspace{\ZSFspaceXXS}%
   \fi
   \begin{itemize}[
     leftmargin=1.2em,
     labelsep=0.4em,
     itemsep=\ZSFprocItemSep,
     topsep=\ZSFspaceXS,
     parsep=0pt,
     partopsep=0pt
   ]}
  {\end{itemize}\par}

% ------------------------------------------------------------
% Goal-System: Bedingung -> Umformung -> Ziel (Inline + Card)
% ------------------------------------------------------------
\newtcolorbox{goalbox}[1][]{
  zsfpanelbox,
  breakable=false,
  colback=white,
  left=\ZSFspaceS,
  right=\ZSFspaceS,
  top=\ZSFtextMinPad,
  bottom=\ZSFtextMinPad,
  title={#1},
  before upper={\parskip=\ZSFspaceS\centering},
  after upper={\par},
}

\newtcbox{\GoalConditionBox}{%
  on line, enhanced,
  boxrule=1.5pt,
  colback=white,
  colframe=\chaptercolorlight!22,
  borderline={0.2pt}{0pt}{black},
  arc=\ZSFcornerInline, outer arc=\ZSFcornerInline,
  left=\ZSFspaceXS, right=\ZSFspaceXS, top=\ZSFspaceXS, bottom=\ZSFspaceXS,
  nobeforeafter,
  tcbox raise base,
}
\newtcbox{\GoalTargetBox}{%
  on line, enhanced,
  boxrule=0.5pt,
  colback=\chaptercolorlight!22,
  colframe=\chaptercolor,
  arc=\ZSFcornerInline, outer arc=\ZSFcornerInline,
  left=\ZSFspaceXS, right=\ZSFspaceXS, top=\ZSFspaceXS, bottom=\ZSFspaceXS,
  nobeforeafter,
  tcbox raise base,
}

\newcommand{\GoalCondition}[1]{\par\centering\GoalConditionBox{\ZSFfontNote #1}\par}
\newcommand{\GoalStep}[1]{\par\centering$\displaystyle #1$\par}
\newcommand{\GoalTarget}[1]{\par\centering\GoalTargetBox{$\displaystyle #1$}\par}
\newcommand{\ZSFDerivationCase}[3]{%
  \GoalCondition{#1}%
  \GoalStep{#2 =}%
  \GoalTarget{#3}%
}
\newcommand{\ZSFDerivationInlineCase}[3]{%
  \GoalCondition{#1}%
  \GoalStep{#2 = \GoalTargetBox{$\displaystyle #3$}}%
}
\newcommand{\GoalCard}[4][]{%
  \begin{goalbox}[#1]
    \GoalCondition{#2}
    #3
    \GoalTarget{#4}
  \end{goalbox}%
}
\newcommand{\GoalPair}[8]{%
  \begin{tcbraster}[raster columns=2, raster column skip=\ZSFspaceM, raster row skip=0pt, raster equal height]
    \GoalCard[#1]{#2}{#3}{#4}
    \GoalCard[#5]{#6}{#7}{#8}
  \end{tcbraster}%
}

% Inline-Danger-Tag (rotes Pill-Label für Stolperfallen im Fliesstext)
\newtcbox{\ZSFdanger}{
  colback=red!80!black, colframe=red!80!black, coltext=white,
  on line, boxsep=0.5pt,
  left=\ZSFspaceS, right=\ZSFspaceS, top=0pt, bottom=0pt,
  boxrule=0pt, arc=\ZSFcornerInline, outer arc=\ZSFcornerInline,
  fontupper=\sffamily\bfseries\ZSFVariableColorsOff
}

% ------------------------------------------------------------
% Valuegrid: kompakte (tabularray-)Wertegrids in einer Panel-Box
% ------------------------------------------------------------
\newtcolorbox{compactgridbox}[1][]{
  zsfpanelbox,
  colframe=black!20,
  breakable=false,
  title={#1},
  boxsep=0pt, left=0pt, right=0pt, top=\ZSFtextMinPad, bottom=\ZSFtextMinPad,
  zsftitlebar,
  boxrule=0.5pt
}
\newcommand{\ZSFvaluegridCommon}{
  colsep=3pt, % MaschKo-Verdichtung: 4 -> 3 (6-Seiten-Ziel)
  rowsep=2pt, % MaschKo-Verdichtung: 3 -> 2
  hlines,
  vlines,
  rows={valign=m}
}
% \begin{valuegrid}{<n_datacols>}[<title>] — Spec wird automatisch generiert.
\NewDocumentEnvironment{valuegrid}{m O{}}
  {\begin{compactgridbox}[#2]\begin{tblr}{colspec={l *{#1}{c}},\ZSFvaluegridCommon}}
  {\end{tblr}\end{compactgridbox}}
% Freie colspec-Variante
\NewDocumentEnvironment{valuegridcustom}{m O{}}
  {\begin{compactgridbox}[#2]\begin{tblr}{colspec={#1},\ZSFvaluegridCommon}}
  {\end{tblr}\end{compactgridbox}}

% Komfort-Wrapper für feste Spaltenzahlen — sparen das Zählen der Datenspalten.
% \begin{valuegridthree}[Titel] ... \end{valuegridthree}  == valuegrid{3}[Titel]
\NewDocumentEnvironment{valuegridtwo}{O{}}{\begin{valuegrid}{2}[#1]}{\end{valuegrid}}
\NewDocumentEnvironment{valuegridthree}{O{}}{\begin{valuegrid}{3}[#1]}{\end{valuegrid}}
\NewDocumentEnvironment{valuegridfour}{O{}}{\begin{valuegrid}{4}[#1]}{\end{valuegrid}}
\NewDocumentEnvironment{valuegridfive}{O{}}{\begin{valuegrid}{5}[#1]}{\end{valuegrid}}
\NewDocumentEnvironment{valuegridsix}{O{}}{\begin{valuegrid}{6}[#1]}{\end{valuegrid}}
\NewDocumentEnvironment{valuegridseven}{O{}}{\begin{valuegrid}{7}[#1]}{\end{valuegrid}}
